% Ban thu: chi giu cac cau co trong CD2

%@ id=TT25.06.003 | trang=1 | so=2 | de=II | vitri=CD2.2,C3.9.5 | muc=8 | tt=0
\begin{bt}
Cho hai biểu thức $P=\dfrac{x+3}{\sqrt{x}-2}$; $Q=\dfrac{\sqrt{x}-1}{\sqrt{x}+2}+\dfrac{5\sqrt{x}-2}{x-4}$, với $x>0$, $x\ne4$.
\begin{cacy}
\item Tính giá trị của biểu thức $P$ khi $x=9$.
\item Rút gọn biểu thức $Q$.
\item Tìm giá trị của $x$ để biểu thức $\dfrac{P}{Q}=4$.
\end{cacy}
\loigiai{
1. Ta có $x=9$ thỏa mãn điều kiện. Thay vào biểu thức $P$ ta có: $P=\dfrac{9+3}{\sqrt{9}-2}=12$. Vậy với $x=9$ thì $P$ có giá trị bằng $12$.\\
2. Ta có:
\[
Q=\dfrac{\sqrt{x}-1}{\sqrt{x}+2}+\dfrac{5\sqrt{x}-2}{(\sqrt{x}+2)(\sqrt{x}-2)}=\dfrac{(\sqrt{x}-1)(\sqrt{x}-2)+5\sqrt{x}-2}{(\sqrt{x}+2)(\sqrt{x}-2)}=\dfrac{x-2\sqrt{x}-\sqrt{x}+2+5\sqrt{x}-2}{(\sqrt{x}+2)(\sqrt{x}-2)}
\]
\[
Q=\dfrac{x+2\sqrt{x}}{(\sqrt{x}+2)(\sqrt{x}-2)}=\dfrac{\sqrt{x}(\sqrt{x}+2)}{(\sqrt{x}+2)(\sqrt{x}-2)}=\dfrac{\sqrt{x}}{\sqrt{x}-2}.
\]
Vậy với $x>0$, $x\ne4$ thì $Q=\dfrac{\sqrt{x}}{\sqrt{x}-2}$.\\
3. Với $x>0$, $x\ne4$ ta có:
\[
\dfrac{P}{Q}=\dfrac{x+3}{\sqrt{x}-2}:\dfrac{\sqrt{x}}{\sqrt{x}-2}=\dfrac{x+3}{\sqrt{x}-2}\cdot\dfrac{\sqrt{x}-2}{\sqrt{x}}=\dfrac{x+3}{\sqrt{x}}.
\]
Để $\dfrac{P}{Q}=4$ thì:
\[
\dfrac{x+3}{\sqrt{x}}=4,\qquad x+3=4\sqrt{x},\qquad x-4\sqrt{x}+3=0,
\]
\[
(\sqrt{x}-1)(\sqrt{x}-3)=0.
\]
TH1: $x=1$ (thỏa mãn).\\
TH2: $x=9$ (thỏa mãn).\\
Vậy với $x=1$ hoặc $x=9$ thì $\dfrac{P}{Q}=4$.
}
\end{bt}
